% Chapitre 6 : la boucle de relecture.
\begin{figure}[tbp]
\centering
\begin{tikzpicture}[x=1cm,y=1cm,
  etape/.style={boite, minimum width=2.9cm, text width=2.75cm, minimum height=1.25cm}]
  \node[etape] (marques) at (0,0) {\textbf{Marques de doute}\\[-1pt]\scriptsize
    \texttt{\textbackslash uncertain}, \texttt{\textbackslash ill} dans la transcription};
  \node[etape] (page) at (4.55,0) {\textbf{Page de relecture}\\[-1pt]\scriptsize
    l'image, la lecture proposée, une question fermée};
  \node[etape] (rep) at (9.1,0) {\textbf{Réponses}\\[-1pt]\scriptsize oui, non,
    incertain, copiées en un bloc};
  \node[etape] (source) at (9.1,-2.4) {\textbf{Appliquées à la source}\\[-1pt]\scriptsize
    le fichier,\\son \mbox{en-tête},\\puis un commit};
  \node[etape, draw=tonregle, fill=tonregle!7] (close) at (4.55,-2.4) {\textbf{Question
    fermée}\\[-1pt]\scriptsize retirée de la page};
  \draw[fleche] (marques) -- node[etiq, above] {fichier JSON} (page);
  \draw[fleche] (page) -- node[etiq, above] {une personne} (rep);
  \draw[fleche] (rep) -- node[etiq, right] {l'assistant} (source);
  \draw[fleche] (source) -- (close);
  \draw[fleche, tonregle] (rep.north) -- ++(0,0.5) -| node[etiq, text=tonregle, pos=0.25,
    above] {incertain : la question est reposée, plus précise} (page.north);
\end{tikzpicture}
\caption{La boucle de relecture. L'assistant marque ses doutes, une personne y répond sur
une page qui montre l'original, et l'assistant applique les réponses dans le fichier
source avant de fermer la question.}
\label{fig:relecture}
\end{figure}
